\begin{frame}{Concepto de análisis financiero}
  \begin{columns}[T,onlytextwidth]
    \column{0.46\textwidth}
    \footnotesize
    El análisis financiero es un proceso técnico de revisión, comparación e interpretación de estados financieros.

    \vspace{0.4em}
    Su finalidad es evaluar posición financiera, desempeño económico y capacidad para generar resultados sostenibles.

    \vspace{0.4em}
    No es solo cálculo: es diagnóstico empresarial.
    \tiempo{0.8 minutos}
    \column{0.50\textwidth}
    \figpng[0.66\textheight]{F01_Proceso_Analisis_Financiero.png}
  \end{columns}
  \note{Definir análisis financiero y destacar interpretación.}
\end{frame}

\begin{frame}{Función del análisis financiero en la empresa}
  \begin{columns}[T,onlytextwidth]
    \column{0.48\textwidth}
    \begin{itemize}
      \item Evaluar desempeño y rentabilidad.
      \item Detectar tendencias y riesgos.
      \item Controlar costos y márgenes.
      \item Comunicar información a usuarios internos.
    \end{itemize}
    \micro{La gerencia usa el análisis financiero para decidir con base en evidencia contable y no solo en percepciones.}
    \tiempo{0.8 minutos}
    \column{0.48\textwidth}
    \figpng[0.62\textheight]{F14_Decisiones_Empresariales.png}
  \end{columns}
  \note{Explicar para qué sirve el análisis financiero en la empresa.}
\end{frame}

\begin{frame}{Ejemplo aplicado: de datos contables a tablero gerencial}
  \begin{columns}[T,onlytextwidth]
    \column{0.42\textwidth}
    \footnotesize
    Una empresa comercial registra diariamente ventas, compras y gastos en su ERP. La gerencia necesita saber si las ventas crecen, si los costos se mantienen controlados y si la utilidad mejora.

    \vspace{0.35em}
    \begin{itemize}
      \item ERP: registra operaciones.
      \item ETL: organiza la información.
      \item Dashboard: presenta alertas.
    \end{itemize}
    \micro{El valor del sistema está en convertir datos en información financiera útil.}
    \tiempo{0.8 minutos}
    \column{0.54\textwidth}
    \flujoERPdashboard
  \end{columns}
  \note{Ejemplo aplicado de ERP a dashboard.}
\end{frame}

\begin{frame}{Estados financieros utilizados}
  \begin{columns}[T,onlytextwidth]
    \column{0.45\textwidth}
    \begin{itemize}
      \item Estado de Situación Financiera.
      \item Estado de Resultados.
      \item Estado de Flujos de Efectivo.
      \item Estado de Cambios en el Patrimonio.
    \end{itemize}
    \micro{Para el caso práctico se priorizan los dos primeros por su aplicación directa al cálculo horizontal y vertical.}
    \tiempo{0.8 minutos}
    \column{0.51\textwidth}
    \figpng[0.62\textheight]{F02_Estado_Situacion_Financiera.png}
  \end{columns}
  \note{Explicar los cuatro estados y priorización del balance y resultados.}
\end{frame}

\begin{frame}{Estado de Situación Financiera}
  \begin{columns}[T,onlytextwidth]
    \column{0.44\textwidth}
    \begin{itemize}
      \item Activos: recursos controlados.
      \item Pasivos: obligaciones.
      \item Patrimonio: participación de propietarios.
    \end{itemize}
    \[
      \text{Activos}=\text{Pasivos}+\text{Patrimonio}
    \]
    \tiempo{0.8 minutos}
    \column{0.52\textwidth}
    \figpng[0.64\textheight]{F02_Estado_Situacion_Financiera.png}
  \end{columns}
  \note{Explicar balance y ecuación contable.}
\end{frame}

\begin{frame}{Ejemplo aplicado: lectura del Estado de Situación Financiera}
  \begin{columns}[T,onlytextwidth]
    \column{0.48\textwidth}
    \footnotesize
    Si la empresa aumenta activos, pero también incrementa deuda, el análisis debe identificar si el crecimiento se financia con recursos propios o con obligaciones.

    \vspace{0.5em}
    Un sistema puede generar alertas cuando el pasivo crece más rápido que el patrimonio o cuando baja la liquidez.
    \tiempo{0.7 minutos}
    \column{0.48\textwidth}
    \figpng[0.62\textheight]{F02_Estado_Situacion_Financiera.png}
  \end{columns}
  \note{Ejemplo aplicado al balance.}
\end{frame}

\begin{frame}{Estado de Resultados}
  \begin{columns}[T,onlytextwidth]
    \column{0.42\textwidth}
    \begin{itemize}
      \item Ingresos.
      \item Costos y gastos.
      \item Impuestos.
      \item Utilidad o pérdida.
    \end{itemize}
    \micro{Permite evaluar si las ventas se convierten en utilidad y si la empresa controla sus costos.}
    \tiempo{0.8 minutos}
    \column{0.54\textwidth}
    \figpng[0.66\textheight]{F03_Estado_Resultados.png}
  \end{columns}
  \note{Explicar estado de resultados.}
\end{frame}

\begin{frame}{Ejemplo aplicado: empresa de servicios digitales}
  \begin{columns}[T,onlytextwidth]
    \column{0.48\textwidth}
    \footnotesize
    Una empresa de software puede vender licencias, soporte y desarrollo. El Estado de Resultados permite evaluar si esos ingresos cubren costos de personal, servidores, licencias y soporte.

    \vspace{0.5em}
    Si los costos técnicos crecen más rápido que las ventas, el margen operativo puede deteriorarse.
    \tiempo{0.7 minutos}
    \column{0.48\textwidth}
    \figpng[0.62\textheight]{F03_Estado_Resultados.png}
  \end{columns}
  \note{Ejemplo de empresa tecnológica.}
\end{frame}

\begin{frame}{Relación entre balance y resultados}
  \begin{columns}[T,onlytextwidth]
    \column{0.46\textwidth}
    \begin{itemize}
      \item Resultados: desempeño del periodo.
      \item Balance: estructura acumulada.
      \item Lectura conjunta: rentabilidad, liquidez y deuda.
    \end{itemize}
    \tiempo{0.7 minutos}
    \column{0.50\textwidth}
    \figpng[0.62\textheight]{F16_Cierre_Conceptual.png}
  \end{columns}
  \note{Explicar relación balance-resultados.}
\end{frame}

\begin{frame}{Ejemplo aplicado: lectura integrada}
  \begin{columns}[T,onlytextwidth]
    \column{0.48\textwidth}
    \footnotesize
    Una empresa puede registrar utilidad positiva en el Estado de Resultados, pero presentar problemas de liquidez si gran parte de sus ventas aún no ha sido cobrada.

    \vspace{0.5em}
    Ejemplo: ventas crecientes, utilidad neta positiva, cuentas por cobrar elevadas y efectivo disponible reducido.

    \vspace{0.5em}
    El análisis debe integrar resultados y situación financiera para evitar conclusiones basadas en un solo reporte.
    \tiempo{0.7 minutos}
    \column{0.48\textwidth}
    \flujoERPdashboard
  \end{columns}
  \note{Ejemplo de lectura integrada.}
\end{frame}
